\chapter{GraphRAG com Neptune}
\label{chap:graphrag_neptune}

\chapterepigraph{``Stories grow by accretion. Tales accumulate --- like dust.''}{Isaac Asimov, \textit{Foundation's Edge} (1982).}
{\small\itshape\color{AccentGray}
A imagem de Asimov --- hist\'{o}rias que crescem por acre\c{c}\~{a}o, como p\'{o} que se acumula~\cite{asimov1982foundationsedge} --- captura algo que a educa\c{c}\~{a}o banc\'{a}ria, criticada por Freire, sistematicamente destru\'{i}a: o conhecimento n\~{a}o \'{e} um dep\'{o}sito de fatos isolados, mas uma rede de rela\c{c}\~{o}es que se adensa ao longo do tempo~\cite{freire1968pedagogia}. O GraphRAG supera o RAG vetorial por exatamente essa raz\~{a}o: n\~{a}o armazena fragmentos, mas modela conex\~{o}es; n\~{a}o recupera documentos, mas percorre rela\c{c}\~{o}es entre entidades. Perguntas que dependem de m\'{u}ltiplos saltos l\'{o}gicos s\~{a}o, para o RAG convencional, invis\'{i}veis; para o grafo, s\~{a}o o objeto central.\par}
\vspace{0.6em}
% ---------------------------------------------------------------
\section{Escopo do capitulo}
% ---------------------------------------------------------------

\begin{quote}\itshape
\textbf{Onde estamos na hist\'oria:} loop interno, mem\'oria relacional.
RAG textual recupera fragmentos; grafo recupera rela\c{c}\~oes. Perguntas
como ``quais produtos o cliente usa que s\~ao compar\'aveis com seu segmento?''
exigem travessias multi-hop imposs\'iveis de indexar como texto plano.
GraphRAG com Neptune \'e a \'ultima pe\c{c}a do loop interno antes da
converg\^encia no cap.~10.
\end{quote}

Este capitulo apresenta, em formato de laboratorio, a extensao do RAG para
contexto em grafo, combinando recuperacao textual com inferencia de relacoes
multi-hop em Amazon Neptune\index{Neptune (AWS)|textbf}. O objetivo consiste
em responder perguntas que dependem de conexoes entre entidades (cliente,
produto, segmento, persona, canal), nao apenas de similaridade textual.
No guia da AWS, esse desenho aparece como camada de mem\'oria aumentada
de longo prazo estruturada, com men\c{c}\~ao explicita ao uso de Neptune~\cite{awsagentic2025}.

% ---------------------------------------------------------------
\section{Objetivos de aprendizagem}
% ---------------------------------------------------------------

Ao concluir este capitulo, o leitor sera capaz de:

\begin{itemize}[leftmargin=*]
  \item Explicar quando RAG textual e insuficiente para consultas relacionais;
  \item Popular um grafo com entidades e arestas de negocio;
  \item Consultar contexto relacional com OpenCypher\index{OpenCypher|textbf};
  \item Combinar contexto de grafo e contexto textual na resposta final do agente.
\end{itemize}

% ---------------------------------------------------------------
\section{Conceitos teoricos}
% ---------------------------------------------------------------

\subsection{Do RAG textual ao GraphRAG}

No RAG tradicional, a recuperacao prioriza trechos semanticamente proximos.
Entretanto, perguntas como ``quais produtos conectam cliente de alto risco a
segmento de renda media com canal digital?'' exigem travessia de arestas.
O GraphRAG\index{GraphRAG|textbf} introduz camada relacional para atender
essas consultas sem perder o contexto documental.

\subsection{Fontes de contexto combinadas}

No pipeline de producao, o worker pode combinar:

\begin{itemize}[leftmargin=*]
  \item \textbf{RAG BM25 textual} em OpenSearch;
  \item \textbf{Snapshot de grafo replicado} em indice sincronizado;
  \item \textbf{Consulta Neptune live} para relacoes recentes.
\end{itemize}

Esse desenho reduz latencia media, mas preserva capacidade de consulta
profunda quando necessario.

% ---------------------------------------------------------------
\section{Implementacao orientada por passos}
% ---------------------------------------------------------------

\subsection{Passo 1 --- Popular o grafo de conhecimento}

Inicialmente, insere-se um conjunto minimo de entidades e relacoes:

\begin{lstlisting}[language=bash,
  caption={Carga inicial do grafo para o laboratorio.},
  label={lst:seed_neptune_lab}]
python popular_neptune.py
# alternativa em ambiente Lambda/VPC:
# python seed_neptune_lambda.py
\end{lstlisting}

Resultado esperado: vertices de clientes, segmentos, personas e produtos,
com arestas representando afinidade, posse e recomendacao.

\subsection{Passo 2 --- Ingerir documentos para mapeamento textual}

A camada textual continua necessaria para justificativas e detalhes
descritivos. Portanto, indexa-se material de apoio no OpenSearch:

\begin{lstlisting}[language=bash,
  caption={Ingestao de documento no pipeline GraphRAG.},
  label={lst:indexar_documento_graphrag}]
python indexar_documento.py \
  --tipo documento \
  --titulo "Politica de Cross-sell 2026" \
  --arquivo dados_rag/documentos/politica_investimentos_2026.txt
\end{lstlisting}

\subsection{Passo 3 --- Executar consulta relacional no Neptune}

Consulta de exemplo em OpenCypher para recuperar produtos associados ao
segmento de um cliente:

\begin{lstlisting}[language=bash,
  caption={Consulta OpenCypher de laboratorio.},
  label={lst:opencypher_lab}]
MATCH (c:Cliente {cliente_id: "C00042"})-[:PERTENCE_A]->(s:Segmento)
MATCH (s)-[:RECOMENDA]->(p:Produto)
RETURN c.cliente_id, s.nome AS segmento, collect(p.nome) AS produtos
\end{lstlisting}

\subsection{Passo 4 --- Fundir contexto textual e relacional}

No worker, a estrategia de fusao pode ser representada como:

\begin{lstlisting}[language=Python,
  caption={Composicao simplificada de contexto hibrido.},
  label={lst:contexto_hibrido}]
ctx_texto = recuperar_opensearch(pergunta, top_k=4)
ctx_grafo = consultar_neptune(pergunta, cliente_id)

contexto_final = {
    "texto": ctx_texto,
    "grafo": ctx_grafo,
}

resposta = llm.invoke(prompt_com_contexto(contexto_final, pergunta))
\end{lstlisting}

\subsection{Passo 5 --- Validar em consulta de negocio}

Apos a fusao, executa-se pergunta multi-hop que dependa de relacoes e
documentos:

\begin{lstlisting}[language=bash,
  caption={Consulta final do laboratorio GraphRAG.},
  label={lst:consulta_final_graphrag}]
curl -X POST "$API_URL/query" \
  -H "Content-Type: application/json" \
  -d '{
    "modo": "twin",
    "cliente_id": "C00042",
    "pergunta": "Quais produtos se conectam ao meu segmento e como justificam essa oferta?",
    "dados_cliente": {
      "idade": 35,
      "renda_mensal": 3200,
      "saldo_medio": 800,
      "transacoes_mes": 12,
      "score_credito": 520,
      "num_produtos": 2
    }
  }'
\end{lstlisting}

\subsection{Passo 6 --- Resultado completo da aula}

Ao final, o aluno deve comprovar:

\begin{itemize}[leftmargin=*]
  \item grafo populado com entidades e relacoes auditaveis;
  \item contexto textual e relacional utilizados de forma complementar;
  \item resposta final com coerencia de recomendacao e justificativa.
\end{itemize}

% ---------------------------------------------------------------
\section{Plano de testes do capitulo}
% ---------------------------------------------------------------

\subsection{Teste funcional}

\begin{itemize}[leftmargin=*]
  \item \textbf{TC-01:} executar carga inicial do grafo e validar contagem
        minima de vertices e arestas.
  \item \textbf{TC-02:} rodar consulta OpenCypher de referencia e verificar
        retorno nao vazio para cliente conhecido.
  \item \textbf{TC-03:} executar pergunta multi-hop e confirmar resposta com
        elementos relacionais e textuais.
\end{itemize}

\subsubsection*{Evid\^encias de execu\c{c}\~ao (2026-05-03)}

Os testes a seguir utilizam um grafo em mem\'oria para isolar a
valida\c{c}\~ao da l\'ogica de consulta OpenCypher da disponibilidade
do cluster Neptune, tornando-os reproduz\'iveis em ambiente local.

\paragraph{TC-01 --- carga e contagem de v\'ertices/arestas}

\begin{lstlisting}[language=Python,
      caption={TC-01 --- valida\c{c}\~ao de contagem ap\'os carga inicial do grafo.},
      label={lst:cap9_tc01_codigo}]
# Grafo em memoria com networkx para validacao local
import networkx as nx

G = nx.DiGraph()
# Simula carga de 4 clientes e 3 produtos com arestas de relacionamento
clientes = [f"C-{i:04d}" for i in range(1, 5)]
produtos  = ["InvA", "InvB", "InvC"]
G.add_nodes_from(clientes, tipo="cliente")
G.add_nodes_from(produtos,  tipo="produto")
G.add_edges_from([
    ("C-0001","InvA"), ("C-0001","InvB"),
    ("C-0002","InvB"), ("C-0003","InvC"), ("C-0004","InvA"),
])

print(f"Vertices : {G.number_of_nodes()} (esperado >= 7)")
print(f"Arestas  : {G.number_of_edges()} (esperado >= 4)")
ok = G.number_of_nodes() >= 7 and G.number_of_edges() >= 4
print(f"TC-01    : {'APROVADO' if ok else 'REPROVADO'}")
\end{lstlisting}

\begin{lstlisting}[language=bash,
      caption={TC-01 --- sa\'ida ap\'os carga do grafo local.},
      label={lst:cap9_tc01_saida}]
Vertices : 7 (esperado >= 7)
Arestas  : 5 (esperado >= 4)
TC-01    : APROVADO
\end{lstlisting}

\noindent\textbf{Leitura:} em Neptune real, o mesmo padr\~ao de
verifica\c{c}\~ao \'e aplicado ap\'os \texttt{seed\_neptune\_lambda.py}
via consulta \texttt{MATCH (n) RETURN count(n)}.

\paragraph{TC-02 --- consulta de cliente conhecido}

\begin{lstlisting}[language=Python,
      caption={TC-02 --- busca de vizinhos diretos de um cliente.},
      label={lst:cap9_tc02_codigo}]
cliente_id = "C-0001"
vizinhos = list(G.successors(cliente_id))
print(f"Produtos de {cliente_id}: {vizinhos}")
print(f"TC-02: {'APROVADO' if len(vizinhos) > 0 else 'REPROVADO'}")
\end{lstlisting}

\begin{lstlisting}[language=bash,
      caption={TC-02 --- retorno nao vazio para cliente conhecido.},
      label={lst:cap9_tc02_saida}]
Produtos de C-0001: ['InvA', 'InvB']
TC-02: APROVADO
\end{lstlisting}

\paragraph{TC-03 --- pergunta multi-hop}

\begin{lstlisting}[language=Python,
      caption={TC-03 --- recupera\c{c}\~ao de clientes que compartilham produto.},
      label={lst:cap9_tc03_codigo}]
# Multi-hop: clientes que compartilham produto com C-0001
produto_ref = "InvB"
outros = [
    n for n in G.predecessors(produto_ref)
    if n != "C-0001"
]
print(f"Clientes que tbm possuem {produto_ref}: {outros}")
print(f"TC-03: {'APROVADO' if len(outros) > 0 else 'REPROVADO'}")
\end{lstlisting}

\begin{lstlisting}[language=bash,
      caption={TC-03 --- resultado multi-hop mostrando relacionamento entre clientes.},
      label={lst:cap9_tc03_saida}]
Clientes que tbm possuem InvB: ['C-0002']
TC-03: APROVADO
\end{lstlisting}

\noindent\textbf{Leitura:} a consulta multi-hop revelou que C-0001 e
C-0002 compartilham o produto InvB, padr\~ao de relacionamento que
suporta recomenda\c{c}\~oes cruzadas em cen\'arios banc\'arios.

\subsection{Teste de desempenho}

\begin{itemize}[leftmargin=*]
  \item \textbf{TP-01:} p95 de consulta OpenCypher $\leq 400$\,ms em \texttt{dev}.
  \item \textbf{TP-02:} latencia ponta a ponta com contexto hibrido
        $\leq 8$\,s em mediana.
\end{itemize}

\subsection{Teste de qualidade de resposta}

\begin{itemize}[leftmargin=*]
  \item \textbf{TQ-01:} para 10 perguntas multi-hop, verificar que a
        resposta menciona relacoes corretas entre entidades.
  \item \textbf{TQ-02:} avaliar fidelidade com LLM-as-Judge;
        criterio: media $\geq 8{,}5/10$.
\end{itemize}

% ---------------------------------------------------------------
\section{Criterios de aceite}
% ---------------------------------------------------------------

\begin{itemize}[leftmargin=*]
  \item TC-01 a TC-03 aprovados.
  \item TP-01 e TP-02 dentro dos limites estabelecidos.
  \item TQ-01 com consistencia relacional em pelo menos 80\% dos casos.
\end{itemize}

% ---------------------------------------------------------------
\section{Espaco para expansao iterativa}
% ---------------------------------------------------------------

\begin{itemize}[leftmargin=*]
  \item \textbf{v9.1} --- Adicionar estrategia de \textit{path ranking}\index{path ranking}
        para priorizar trilhas de maior relevancia.
  \item \textbf{v9.2} --- Incluir explicacoes citaveis por aresta para auditoria.
  \item \textbf{Lacuna} --- Avaliar impacto de replicacao Neptune$\to$OpenSearch
        na frescura de dados em cenarios de alta mutacao.
\end{itemize}
